\documentclass[10pt,a4paper]{amsart}
\usepackage[margin=19mm]{geometry}
\usepackage{amsmath,amssymb,mathtools}
\usepackage[scr=rsfs,cal=euler]{mathalfa}
\usepackage{xfrac}
\usepackage[unicode,hidelinks,bookmarks=false]{hyperref}
\newcommand{\ie}{i.e.\ }
\newcommand{\cf}{cf.\ }
\newcommand{\resp}{respectively\ }
\newcommand{\Subsection}{\subsection}
\providecommand{\nobreakdash}{\nobreak\hbox{-}}
\setlength{\emergencystretch}{2em}
\multlinegap0pt
\usepackage{amssymb}
\usepackage{xcolor}
\newtheorem{theorem}{Theorem}
\title{A Geometric Profile of Semantic Information in Text:\\ Frame-Conditional Uniqueness and a Trade-Off Triangle for Scalar Summaries}
\author{Dmitriy Kompaneets}
\date{Source: arXiv:2606.11222v1 (CC BY 4.0)}
\begin{document}
\maketitle

\noindent\emph{Source and licence.} A Geometric Profile of Semantic Information in Text:\\ Frame-Conditional Uniqueness and a Trade-Off Triangle for Scalar Summaries. Version arXiv:2606.11222v1 (CC BY 4.0), distributed by arXiv under the Creative Commons Attribution 4.0 International licence, http://creativecommons.org/licenses/by/4.0/, as recorded in the arXiv licence metadata for this exact version and shown on the abstract page https://arxiv.org/abs/2606.11222v1. Full licence notice: \texttt{licenses/arxiv-2606.11222-CC-BY-4.0.txt}.\par\smallskip

\noindent\textit{Editorial scope.} Counted unit: the article's theorem thm:dpp-breadth (source0.tex lines 339--351). The article's complete original proof of it is delivered with it (source0.tex lines 353--365, one \texttt{proof} environment of 13 lines). No mathematics is written, repaired or re-derived: the statement and the proof are the article's own lines, in the article's own notation, and no retained line is substituted or re-flowed.  No further source-local context is needed: every cross-reference of every retained line resolves inside the two delivered spans.  Nothing was omitted from any retained span, so the package carries no omission record.  No retained line contains a citation, so the wrapper carries no bibliography and no bibliography entry.  No retained line contains a figure environment, an image-inclusion command or drawing code, so no image is delivered and the package declares no asset dependency.  Editorial formatting only, in addition: an AMS \texttt{amsart} preamble that restates the article's own declarations and macros used by the retained lines, namely the environments \texttt{theorem} on separate counters, together with the standard packages those lines need.  The retained environments print in this document as Theorem 1 in order of appearance, whereas the article prints the same items with the numbering of its own full document.  The complete native source of the article as distributed in the arXiv e-print for this exact version is bundled unchanged as \texttt{sources/202530/source0.tex}.
\begin{theorem}[DPP--breadth structural decomposition]\label{thm:dpp-breadth}
Let $G_S = [\langle x_i, x_j \rangle]_{i,j \in S}$ be the Gram matrix of cosines restricted to $S$. Then
\[
\log \det L_S \;=\; 2 \sum_{i \in S} \log q_i \;+\; \log \det G_S.
\]
Moreover, let $\{\lambda_k\}_{k=1}^{|S|}$ denote the eigenvalues of $G_S$, let $D_{\mathrm{eff}}(G_S) = \exp(H(\hat{\lambda}))$ with $\hat{\lambda}_k = \lambda_k / \sum_j \lambda_j$, and let $R(S) = 1 - |S|^{-1}\sum_{i \in S} \langle x_i, \mu_S \rangle / \|\mu_S\|$ where $\mu_S = |S|^{-1} \sum_{i \in S} x_i$. Under the \emph{near-isotropic} regime in which the pairwise cosines satisfy $|\langle x_i, x_j \rangle - c| \leq \varepsilon$ for some $c \in [0, 1)$ and all $i \neq j \in S^*$, the following hold:
\begin{enumerate}
\item[(i)] $\log \det G_S = \sum_k \log \lambda_k$, and the spectrum of $G_S$ concentrates around $(1-c)$ with one eigenvalue at $1 + (|S|-1)c$.
\item[(ii)] $D_{\mathrm{eff}}(G_S) = |S| \cdot \big(1 + O(\varepsilon) + O(|S| c^2)\big)$.
\item[(iii)] $R(S) = 1 - \sqrt{(1 + (|S|-1)c)/|S|} \cdot |S|^{-1} + O(\varepsilon)$, monotone increasing in $|S|$ and decreasing in $c$.
\item[(iv)] Setting $B_S := D_{\mathrm{eff}}(G_S) \cdot R(S)$, there exists a universal function $h: \mathbb{R}_{\geq 0}^2 \to \mathbb{R}$, $h(|S|, c)$, such that $\log \det G_S = h(|S|, c) \cdot (1 + O(\varepsilon))$ and $B_S = h'(|S|, c) \cdot (1 + O(\varepsilon))$, with both $h, h'$ strictly increasing in $|S|$ and strictly decreasing in $c$.
\end{enumerate}
\end{theorem}

\begin{proof}[Proof sketch]
The first identity follows from the multiplicative property of determinants:
\[
\log \det L_S \;=\; \log \det\!\big(\mathrm{diag}(q_S)\, G_S\, \mathrm{diag}(q_S)\big) \;=\; 2 \sum_{i \in S} \log q_i + \log \det G_S.
\]
Under the uniform-correlation model $G_S = (1-c)\,I + c\,J$ (where $J$ is all-ones), the spectrum is one eigenvalue equal to $1 + (|S|-1)c$ and $|S|-1$ eigenvalues equal to $1-c$. The $\varepsilon$ perturbation is controlled by Weyl's inequality. Parts (i)--(iii) follow by direct substitution.

For (iv), both $\log \det G_S$ and $B_S$ are smooth functions of $(|S|, c)$ through the spectrum. In particular,
\[
\log \det G_S = \log\!\big(1 + (|S|-1)c\big) + (|S|-1)\log(1-c),
\]
which is strictly decreasing in $c$ and strictly increasing in $|S|$ for $c \in [0,1)$. $D_{\mathrm{eff}} \cdot R$ inherits the same monotonicities from (ii) and (iii). Both are therefore co-monotone in $(|S|, c)$; the empirical $\rho = 0.985$ ($n = 507$ chapters) realizes this co-monotonicity under the approximate near-isotropy that natural-text embeddings satisfy.
\end{proof}

\end{document}
